\section*{第 \ref{ch:b1:arith} 章 --- 整数算术}

\begin{solution}{exo:b1:arith:1}
$1001 = 1 \times 777 + 224$; $777 = 3 \times 224 + 105$; $224 = 2
\times 105 + 14$；$105 = 7 \times 14 + 7$；$14 = 2 \times 7 + 0$。由此
$\gcd(1001, 777) = 7$。回代得：
\[
7 = 105 - 7 \times 14
= 105 - 7(224 - 2\times 105) = 15 \times 105 - 7 \times 224
\]
\[
= 15(777 - 3\times 224) - 7\times 224 = 15 \times 777 - 52 \times 224
= 15 \times 777 - 52(1001 - 777) = 67 \times 777 - 52 \times 1001 .
\]
检验：$67 \times 777 = 52\,059$ 与 $52 \times 1001 = 52\,052$；
差为 $7$。贝祖对：$(u, v) = (-52, 67)$，使 $1001u + 777v =
7$。
\end{solution}

\begin{solution}{exo:b1:arith:2}
由 $10 \equiv 1 \pmod 9$，得 $10^k \equiv 1$，故 $\sum_k d_k 10^k
\equiv \sum_k d_k \pmod 9$。由 $10 \equiv -1 \pmod{11}$，得 $10^k
\equiv (-1)^k$，故该整数与交错和
$\sum_k (-1)^k d_k$ 模 $11$ 同余（自\emph{个位}数字起
符号为 $+$）。

$123\,456\,789$：数字和 $45 \equiv 0 \pmod 9$。交错和
自个位起：$9 - 8 + 7 - 6 + 5 - 4 + 3 - 2 + 1 = 5$，故该数
$\equiv 5 \pmod{11}$。
\end{solution}

\begin{solution}{exo:b1:arith:3}
欧几里得算法：$237 = 2 \times 91 + 55$；$91 = 1 \times 55 + 36$；$55 = 1
\times 36 + 19$；$36 = 1 \times 19 + 17$；$19 = 1 \times 17 + 2$；$17
= 8 \times 2 + 1$。回代得：
\[
1 = 17 - 8\times 2 = 17 - 8(19 - 17) = 9\times 17 - 8\times 19
= 9(36 - 19) - 8\times 19 = 9\times 36 - 17\times 19
\]
\[
= 9\times 36 - 17(55 - 36) = 26\times 36 - 17\times 55
= 26(91 - 55) - 17\times 55 = 26\times 91 - 43\times 55
\]
\[
= 26\times 91 - 43(237 - 2\times 91) = 112 \times 91 - 43 \times 237.
\]
故 $91 \times 112 \equiv 1 \pmod{237}$，解为 $x \equiv
112 \pmod{237}$。（检验：$91 \times 112 = 10\,192 = 43 \times 237 +
1$。）
\end{solution}

\begin{solution}{exo:b1:arith:4}
$\gcd(17, 39) = 1$：欧几里得算法给出 $39 = 2\times 17 + 5$，$17 = 3\times
5 + 2$，$5 = 2\times 2 + 1$，回代得
\[
1 = 5 - 2\times 2 = 5 - 2(17 - 3\times 5) = 7\times 5 - 2\times 17
= 7(39 - 2\times 17) - 2\times 17 = 7\times 39 - 16\times 17 .
\]
特解为 $(x_0, y_0) = (-16, 7)$。齐次方程
$17x + 39y = 0$ 的通解：$x = 39k$，$y = -17k$（由于
$17 \mid 39y$ 且 $\gcd(17,39) = 1$ 推出 $17 \mid y$ ---
高斯引理）。于是
\[
(x, y) = (-16 + 39k,\; 7 - 17k), \qquad k \in \Z .
\]
右端为 $5$ 时，将特解乘以 $5$：
$(x, y) = (-80 + 39k,\; 35 - 17k)$，$k \in \Z$。
\end{solution}

\begin{solution}{exo:b1:arith:5}
对每个质数 $p$，记 $\alpha = v_p(a)$，$\beta = v_p(b)$，则
\[
v_p\bigl(\gcd(a,b)\bigr) + v_p\bigl(\operatorname{lcm}(a,b)\bigr)
= \min(\alpha, \beta) + \max(\alpha, \beta)
= \alpha + \beta = v_p(ab) .
\]
两个正整数若在每个质数处赋值相同，则必相等
（见\cref{prop:b1:arith:valuation}），故 $\gcd(a,b)\operatorname{lcm}(a,b)
= ab$。
\end{solution}

\begin{solution}{exo:b1:arith:6}
赋值计算：$\gcd(a, b) = 2^{\min(10,6)} 3^{\min(4,7)} 5^{\min(2,0)}
7^{\min(0,1)} = 2^6\, 3^4 = 5184$;
$\operatorname{lcm}(a,b) = 2^{10}\, 3^7\, 5^2\, 7$.

因数个数：$n = \prod p_i^{\alpha_i}$ 的正因数恰为
形如 $\prod p_i^{\beta_i}$ 的数，其中 $0 \leq \beta_i \leq
\alpha_i$（见\cref{prop:b1:arith:valuation}）；这些取法
相互独立，故共有 $\prod_i (\alpha_i + 1)$ 个因数。对 $a$：
$(10+1)(4+1)(2+1) = 165$。
\end{solution}

\begin{solution}{exo:b1:arith:7}
设 $\sqrt p = \frac rq$，其中 $r, q \in \N^*$，即 $p q^2 =
r^2$。两边取 $v_p$：$v_p(pq^2) = 1 + 2v_p(q)$ 为奇数，而 $v_p(r^2)
= 2 v_p(r)$ 为偶数。一个整数不可能同时具有奇、偶两种
$p$ 赋值：矛盾。故 $\sqrt p \notin \Q$。
\end{solution}

\begin{solution}{exo:b1:arith:8}
\emph{一般事实。}由 $\gcd(m,n) = 1$，贝祖等式给出 $mu + nv = 1$。
取 $x_0 = b\,mu + a\,nv$。则 $x_0 \equiv a\,nv \equiv a(1 - mu)
\equiv a \pmod m$，同理 $x_0 \equiv b \pmod n$，存在性得证。若
$x$ 与 $x'$ 为两个解，则 $m$ 与 $n$ 都整除 $x - x'$，
故 $mn \mid x - x'$（见\cref{thm:b1:arith:gauss}~(2)），于是在模
$mn$ 下唯一。

\emph{数值计算：}$m = 7$，$n = 11$：$7 \times (-3) + 11 \times 2 =
1$。故 $x_0 = 5 \times 7 \times (-3) + 2 \times 11 \times 2 = -105 +
44 = -61 \equiv 16 \pmod{77}$。检验：$16 = 2\times 7 + 2 \equiv 2
\pmod 7$；$16 = 11 + 5 \equiv 5 \pmod{11}$。所求解为 $x \equiv 16
\pmod{77}$。
\end{solution}

\begin{solution}{exo:b1:arith:9}
模 $7$：由费马定理得 $3^6 \equiv 1$，且 $1000 = 6 \times 166 + 4$，
故 $3^{1000} \equiv 3^4 = 81 \equiv 4 \pmod 7$。

求 $7^{100}$ 的末两位：分别取模 $4$ 与模 $25$。模 $4$ 时，$7
\equiv -1$，故 $7^{100} \equiv 1$。模 $25$ 时，$7^2 = 49 \equiv -1$，于是
$7^4 \equiv 1$，且 $7^{100} = (7^4)^{25} \equiv 1$。由中国剩余定理
（\cref{exo:b1:arith:8}），$7^{100} \equiv 1
\pmod{100}$，故末两位为 $01$。
\end{solution}

\begin{solution}{exo:b1:arith:10}
设 $m = nq + r$，$0 \leq r < n$。于是
\[
2^m - 1 = 2^r\bigl(2^{nq} - 1\bigr) + 2^r - 1,
\]
且 $2^n - 1$ 整除 $2^{nq} - 1 = (2^n - 1)(2^{n(q-1)} + \dots +
1)$。于是模 $2^n - 1$ 有 $\;2^m - 1 \equiv 2^r - 1$，且由 $0 \leq
2^r - 1 < 2^n - 1$，这\emph{正是}带余除法中的余数。

因此对数对 $(2^m - 1, 2^n - 1)$ 所作的欧几里得算法
与对 $(m, n)$ 所作的算法逐步对应：每次
除法在上方将 $(m, n)$ 换为 $(n, r)$，在下方将 $(2^m - 1,
2^n - 1)$ 换为 $(2^n - 1, 2^r - 1)$。上方的算法
在 $\gcd(m,n)$ 处终止，故下方的算法在
$2^{\gcd(m,n)} - 1$ 处终止。
\end{solution}

\begin{solution}{exo:b1:arith:11}
先解 $x^2 \equiv 1 \pmod p$：$p \mid (x-1)(x+1)$，由
欧几里得引理得 $x \equiv 1$ 或 $x \equiv -1 \pmod p$。

在乘积 $(p-1)! = 1 \times 2 \times \dots \times (p-1)$ 中，每个
因子 $a$ 模 $p$ 可逆，其逆元 $a^{-1}$ 也
是其中的因子（\cref{prop:b1:arith:invmod}）。将每个 $a$ 与
$a^{-1}$：各对之积为 $1$，但自配对的因子
（$a = a^{-1}$，即 $a^2 \equiv 1$）单独保留，这些因子
恰为 $1$ 与 $p - 1$。于是
\[
(p-1)! \equiv 1 \times (p - 1) \equiv -1 \pmod p .
\]
（当 $p = 2$ 时：$1! = 1 \equiv -1 \pmod 2$；
配对论证退化，但结论仍然成立。）

\emph{逆命题。}设 $n \geq 2$ 为合数，$n = ab$，其中 $1 < a
\leq b < n$。若 $a < b$，则二者都是 $(n-1)!$ 的不同因子，
故 $n \mid (n-1)!$，且 $(n-1)! \equiv 0 \not\equiv -1$。若 $a = b$
（即 $n = a^2$）：当 $a \geq 3$ 时，$a$ 与 $2a$ 均 $< n$，于是
$n = a^2 \mid a \times 2a \mid (n-1)!$，结论相同；当 $n = 4$ 时，
$(n-1)! = 6 \equiv 2 \not\equiv -1 \pmod 4$。
\end{solution}

\begin{solution}{exo:b1:arith:12}
\begin{enumerate}
  \item 归纳法。当 $n = 1$ 时：$F_0 = 3 = F_1 - 2 = 5 - 2$。
        假设 $F_0\cdots F_{n-1} = F_n - 2$：
        \[
        F_0 \cdots F_n = (F_n - 2)F_n
        = \bigl(2^{2^n} - 1\bigr)\bigl(2^{2^n} + 1\bigr)
        = 2^{2^{n+1}} - 1 = F_{n+1} - 2 .
        \]
  \item 设 $m < n$，$d = \gcd(F_m, F_n)$。由 (1)，$F_m$
        整除 $F_n - 2$，故 $d$ 同时整除 $F_n$ 与 $F_n -
        2$，从而整除 $2$。但每个费马数都是奇数，故
        $d = 1$。
  \item 每个 $F_n \geq 3$ 都有一个质因数 $p_n$
        （\cref{thm:b1:arith:euclidprimes}的第一步）。若 $m
        \neq n$，则 $p_m \neq p_n$，因为公共质因数
        会整除 $\gcd(F_m, F_n) = 1$。映射 $n \mapsto p_n$
        因而是从 $\N$ 到质数集的单射：质数
        有无穷多个。
\end{enumerate}
\end{solution}

\begin{solution}{pb:b1:arith:1}
\textbf{1.} $10! = 3\,628\,800$，末尾有两个零。逐因子计算赋值：
$2$ 的幂来自 $2, 4 = 2^2, 6, 8 = 2^3,
10$，合计 $v_2(10!) = 1 + 2 + 1 + 3 + 1 = 8$；$5$ 的幂来自
$5$ 与 $10$：$v_5(10!) = 2$。末尾零的个数 $=
\min(v_2, v_5) = 2$，二者一致。

\textbf{2.} 写出带余除法 $\lfloor x\rfloor = nq +
r$，$0 \leq r \leq n - 1$。于是 $x = nq + r + \{x\}$，其中 $0 \leq
r + \{x\} < n$，故 $\lfloor x/n \rfloor = q = \bigl\lfloor \lfloor
x \rfloor / n \bigr\rfloor$。

\textbf{3.} 设 $\alpha = v_p(a) \leq \beta = v_p(b)$（必要时可交换
），并记 $a = p^\alpha a'$，$b = p^\beta b'$，其中 $p
\nmid a', b'$。于是 $a + b = p^\alpha\bigl(a' + p^{\beta -
\alpha}b'\bigr)$，故 $v_p(a + b) \geq \alpha = \min$。若 $\alpha <
\beta$，则括号内为 $a' + p^{\beta-\alpha}b' \equiv a'
\not\equiv 0 \pmod p$，故赋值恰为 $\alpha$。

\textbf{4.} $m$ 的倍数在 $\intint1n$ 中为 $m, 2m, \dots,
qm$ 其中 $q$ 为满足 $qm \leq n$ 的最大整数，即 $q =
\lfloor n/m \rfloor$。

\textbf{5.} 由唯一分解定理，$v_p(n!) = \sum_{j=1}^{n}
v_p(j)$。换一种方式计数：每个 $j$ 贡献 $v_p(j) =
\#\{k \geq 1 : p^k \mid j\}$，于是
\[
v_p(n!) = \sum_{j=1}^n \#\{k : p^k \mid j\}
= \sum_{k\geq1} \#\{j \leq n : p^k \mid j\}
= \sum_{k\geq1} \Bigl\lfloor \frac n{p^k} \Bigr\rfloor
\]
由第 4 题 --- 勒让德公式。和式为有限和：
$p^k > n$ 的项均为零。

\textbf{6.} $v_5(1000!) = 200 + 40 + 8 + 1 = 249$（依次除以
$5, 25, 125, 625$）；$v_2(1000!) = 500 + 250 + 125 + 62 + 31 + 15 +
7 + 3 + 1 = 994$。$1000!$ 末尾的零：每个零消耗
一个 $2$ 与一个 $5$，故共有 $\min(994, 249) = 249$ 个零。

\textbf{7.} 设 $n = \sum_i a_ip^i$，由第 2 题得 $\lfloor
n/p^k \rfloor = \sum_{i \geq k} a_ip^{i-k}$（截断 $p$ 进制
展开）。对 $k \geq 1$ 求和，并交换两个有限
和的次序：
\[
v_p(n!) = \sum_{i\geq1} a_i \sum_{k=1}^{i} p^{i-k}
= \sum_{i\geq0} a_i\,\frac{p^i - 1}{p - 1}
= \frac{n - s_p(n)}{p - 1} .
\]

\textbf{8.} 取 $p = 2$，则 $v_2(n!) = n - s_2(n)$。由 $n \geq 1$
有 $s_2(n) \geq 1$，故总有 $v_2(n!) \leq n - 1 < n$，即 $2^n \nmid
n!$。又 $v_2(n!) = n - 1$ 当且仅当 $s_2(n) = 1$，当且仅当 $n$ 是
$2$ 的幂。

\textbf{9.} $n$ 恰有 $\lfloor \log_p n \rfloor + 1$ 位 $p$ 进制数字，
每位数字至多为 $p - 1$，故 $1 \leq s_p(n) \leq
(p-1)\bigl(\log_p(n) + 1\bigr)$。代入第 7 题：
\[
\frac n{p-1} - \log_p(n) - 1 \;\leq\; v_p(n!) \;<\; \frac n{p-1},
\]
两边除以 $n$，得 $\frac{v_p(n!)}n \to \frac1{p-1}$。

\textbf{10.} $Z(n) - Z(n-1) = v_5(n!/(n-1)!) = v_5(n)$：末尾零的
个数跳增 $v_5(n)$（仅在 $5$ 的倍数处），
其余时刻保持不变。$Z(24) = \lfloor24/5\rfloor = 4$，$Z(25) =
5 + 1 = 6$：当 $n = 25$ 时，计数由 $4$ 直接跳至 $6$
（$v_5(25) = 2$），由于 $Z$ 单调不减，且在此之前 $Z \leq 4$，
在此之后 $Z \geq 6$，故取值 $5$ 永不出现：
不存在恰好以五个零结尾的阶乘。

\textbf{11.} 记 $x = \lfloor x\rfloor + \{x\}$，则 $\lfloor x +
y\rfloor = \lfloor x\rfloor + \lfloor y\rfloor + \lfloor \{x\} +
\{y\}\rfloor$，且 $0 \leq \{x\} + \{y\} < 2$ 使末项的取整
为 $0$ 或 $1$。于是三次应用勒让德公式，
\[
v_p\binom{m+n}m = v_p\bigl((m{+}n)!\bigr) - v_p(m!) - v_p(n!)
= \sum_{k\geq1}\Bigl(
\Bigl\lfloor\frac{m+n}{p^k}\Bigr\rfloor
- \Bigl\lfloor\frac{m}{p^k}\Bigr\rfloor
- \Bigl\lfloor\frac{n}{p^k}\Bigr\rfloor\Bigr),
\]
该和是由 $0$ 与 $1$ 构成的有限和（对 $x =
m/p^k$，$y = n/p^k$）。

\textbf{12.} 固定 $k \geq 1$，写出 $m = p^km_1 + m_0$，$n =
p^kn_1 + n_0$，其中 $0 \leq m_0, n_0 < p^k$（欧几里得除法：
$m_0$ 即由低 $k$ 位数字（取自 $m$）组成的数）。于是
\[
\Bigl\lfloor\frac{m+n}{p^k}\Bigr\rfloor
- \Bigl\lfloor\frac m{p^k}\Bigr\rfloor
- \Bigl\lfloor\frac n{p^k}\Bigr\rfloor
= \Bigl\lfloor\frac{m_0 + n_0}{p^k}\Bigr\rfloor ,
\]
其值为 $1$（当 $m_0 + n_0 \geq p^k$ 时）；否则为 $0$。但 $m_0 +
n_0 \geq p^k$ 恰说明：低 $k$ 位数字（取自 $m$）
与 $n$ 的对应低位相加溢出到第 $k$ 位 --- 即竖式加法中进入第 $k$ 位
的进位。对 $k$ 求和：
$v_p\binom{m+n}m$ 即 $p$ 进制加法
$m + n$ 中的进位个数。（库默尔，1852 年。）

\textbf{13.} 对 $m = j$、$n = p^k - j$ 应用库默尔定理，其和为 $p^k =
(1\underbrace{0\cdots0}_{k})_p$。记 $a = v_p(j)$，则 $p$ 进制下
$j$ 的各位数字在第 $0, \dots, a-1$ 位上为 $0$，
而第 $a$ 位上的数字非零。$p^k - j$ 的各位数字在
第 $a$ 位以下也为 $0$（$p^k - j = p^a(p^{k-a} - j/p^a)$）。
在第 $a$ 位上，两个非零数字之和须为 $p$（结果数字为 $0$）：
故有一次进位；在第 $a+1, \dots, k-1$ 位上，数字与
进入的进位之和为 $p$（结果数字仍为 $0$），进位
向高位传递。共 $k - a$ 次进位，故 $v_p\binom{p^k}j = k -
v_p(j)$。当 $k = 1$ 时，$v_p\binom pj = 1$（$0 < j < p$），
即 \cref{thm:b1:arith:fermat} 所用的整除性。

\textbf{14.} 由第 7 题的数字形式，利用 $s_2(2n) =
s_2(n)$（末尾补一个零数字）：
\[
v_2\binom{2n}n = \bigl(2n - s_2(2n)\bigr) - 2\bigl(n -
s_2(n)\bigr) = 2s_2(n) - s_2(2n) = s_2(n) \geq 1 :
\]
$\binom{2n}n$ 总为偶数，且 $v_2 = 1$（即 $\binom{2n}n
\equiv 2 \pmod 4$）当且仅当 $s_2(n) = 1$，亦即 $n$ 为
$2$ 的幂。

\textbf{15.} 由范德蒙德恒等式（取 $m = n = k = p$）：$\binom{2p}p =
\sum_{j=0}^p \binom pj\binom p{p-j} = \sum_{j=0}^p \binom pj^2$.
当 $0 < j < p$ 时，由第 13 题得 $p \mid \binom pj$，故 $\binom
pj^2 \equiv 0 \pmod p$；两端项给出 $1 + 1$：
$\binom{2p}p \equiv 2 \pmod p$。

\textbf{16.} 在 $3$ 进制下：$500 = 486 + 9 + 3 + 2$，数字（由低位
到高位）$(2, 1, 1, 0, 0, 2)$，故 $s_3(500) = 6$；又 $1000 = 729 +
243 + 27 + 1$，数字为 $(1, 0, 0, 1, 0, 1, 1)$，故 $s_3(1000) = 4$。
\emph{库默尔：}计算 $500 + 500$（在 $3$ 进制下）。第 $0$ 位：$2 + 2 =
4$，数字为 $1$，进位 $1$；第 $1$ 位：$1 + 1 + 1 = 3$，数字为 $0$
进位 $1$；第 $2$ 位：$1 + 1 + 1 = 3$，数字为 $0$，进位 $1$；
第 $3$ 位：$0 + 0 + 1 = 1$，无进位；第 $4$ 位：$0$；
第 $5$ 位：$2 + 2 = 4$，数字为 $1$，进位 $1$；第 $6$ 位：
进位落入此位，数字为 $1$。共四次进位，即 $v_3\binom{1000}{500} = 4$。
\emph{勒让德：}$v_3(1000!) = \frac{1000 - 4}2 = 498$，而
$v_3(500!) = \frac{500 - 6}2 = 247$，故 $v_3\binom{1000}{500} =
498 - 2\times247 = 4$。两种算法结果一致 --- 且
加法所得数字 $(1, 0, 0, 1, 0, 1, 1)$ 还原出 $1000$，这正是
应有的结果。

\textbf{17.} 由库默尔定理（取 $p = 2$，$m = k$，$n' = n - k$）：
$\binom nk$ 为奇数当且仅当加法 $k + (n - k)$ 在 $2$ 进制下
无进位，即在每一位上，数字满足 $k_i + (n -
k)_i = n_i$；此时 $k_i \leq n_i$ 对所有 $i$ 成立。反之，
若 $k_i \leq n_i$ 对所有 $i$ 成立，则以数字 $n_i -
k_i$ 所表示的数恰为 $n - k$，且该加法无进位。同理可证
在 $p$ 进制下：$p \nmid \binom nk$ 当且仅当 $p$ 进制下 $k$ 的每位数字
都不超过 $n$ 中相应的数字。

\textbf{18.} 计数 $k \in \intint0n$ 中各位满足
$k_i \leq n_i$ 的数：$k$ 的每一位可在
$n_i + 1$ 个值中独立选取，共 $\prod_i (n_i + 1)$ 种选法；取基 $2$ 时
此数为 $2^{\#\{i : n_i = 1\}} = 2^{s_2(n)}$。第 $4 =
(100)_2$ 行：$2^1 = 2$ 个奇数项 --- 事实上 $1, 4, 6, 4, 1$ 中
奇数项只在两端。第 $5 = (101)_2$ 行：$2^2 = 4$ --- 事实上
$1, 5, 10, 10, 5, 1$。

\textbf{19.} 行中内部各项皆为偶数 $\iff$ 该行恰有
$2$ 个奇数项（两端项总是奇数）$\iff 2^{s_2(n)} =
2 \iff s_2(n) = 1 \iff n$ 是 $2$ 的幂。

\textbf{20.} 在第 11 题的和式中，第 $k$ 项
一旦 $p^k > m + n$ 即为零（此时三个取整函数皆相等，事实上
首项为 $0$，当 $p^k > m+n$ 时；更简单地说，每一项都为 $0$）。
因此非零项至多有 $\lfloor \log_p(m+n)\rfloor$ 个，每项
之值为 $1$：$a = v_p\binom{m+n}m \leq \log_p(m+n)$，即 $p^a \leq
m + n$。

\textbf{21.} 对每个质数 $p$，$v_p\bigl(\operatorname{lcm}(1,
\dots, 2n)\bigr) = \lfloor\log_p(2n)\rfloor$（
$p$ 的不超过 $2n$ 的最大幂出现在 $1, \dots, 2n$ 中）。
取第 20 题中 $m = n$，得 $v_p\binom{2n}n \leq \lfloor\log_p(2n)\rfloor$
对每个 $p$ 成立。由\cref{prop:b1:arith:valuation}可知，$\binom{2n}n
\mid \operatorname{lcm}(1, \dots, 2n)$。关于大小，比值
$\binom{2n}{k+1}/\binom{2n}k = \frac{2n-k}{k+1} \geq 1$ 当且仅当
$k < n$ 时成立，故中央项是 $2n + 1$ 个
项中最大者，即第 $2n$ 行，于是 $4^n = \sum_k \binom{2n}k \leq
(2n+1)\binom{2n}n$。综合得：
\[
\operatorname{lcm}(1, \dots, 2n) \geq \binom{2n}n \geq
\frac{4^n}{2n + 1} .
\]
若小于 $2n$ 的质数很少，则最小公倍数不可能如此
大：最小公倍数的指数增长，是一种定量标志，反映
质数的丰富程度。

\textbf{22.} $Z(n) = \sum_k\lfloor n/5^k\rfloor \approx \frac
n4$，故取 $n = 4 \times 2026 = 8104$ 附近：$Z(8104) = 1620 +
324 + 64 + 12 + 2 = 2022$。以 $5$ 的倍数逐步增大：$Z(8110)
= 2024$，$Z(8115) = 2025$，且
\[
Z(8120) = 1624 + 324 + 64 + 12 + 2 = 2026 .
\]
由于 $Z$ 在 $5$ 的相邻倍数之间为常值，且 $Z(8119) =
Z(8115) = 2025$，于是最小的 $n$ 使末尾零至少有
$2026$ 个，即 $n = 8120$。

\textbf{23.} 在基 $7$ 下：$50 = 49 + 1$，各位数字（由低到高）为 $(1, 0,
1)$。计算 $50 + 50$：第 $0$ 位 $1 + 1 = 2 < 7$，无进位；
第 $1$ 位 $0 + 0 = 0$；第 $2$ 位 $1 + 1 = 2 < 7$，
亦无进位。由此，库默尔定理给出 $v_7\binom{100}{50} = 0$，即 $7 \nmid
\binom{100}{50}$。由勒让德公式亦可验证：$v_7(100!) = \lfloor 100/7
\rfloor + \lfloor 100/49 \rfloor = 14 + 2 = 16$，且 $v_7(50!) = 7
+ 1 = 8$，故 $v_7\binom{100}{50} = 16 - 2\times8 = 0$。

\textbf{24.} (i) 唯一分解是
$v_p$ 的定义及其可加性之基础，由此得勒让德公式
及每一条整除性结论
（\cref{prop:b1:arith:valuation}）。(ii) 带余除法产生了
第 2 题的截断恒等式，以及分解 $m = p^km_1 +
m_0$ 以隔离进位（第 12 题）。(iii) 计数方面：
$m$ 的倍数计数（第 4 题）、各位选法之积
（第 18 题），以及行和界 $4^n \leq
(2n+1)\binom{2n}n$（第 21 题）均
属\cref{ch:b1:counting}中的计数论证。

\textbf{25.} 勒让德公式把“$p$ 的几次幂整除
$n!$”化为基 $p$ 下的数字运算；库默尔则把二项式系数的答案
压缩为一次加法中的进位；
整除性这一看似关于巨大数的整体性质，
实则可逐位读出的局部性质。第 16 题便是
范例：经手算得到四次进位，便可确定
一个有数百位的数中 $3$ 的确切幂次。而第 21 题
则显示同一组思想已触及深层问题：一个
关于 $\operatorname{lcm}(1, \dots, 2n)$ 的指数型下界，是
迈向质数定理的首个完全初等的一步，
其证明远超本卷范围。整套工具——带余除法、最大公因数、赋值——
在多项式上被重新演练，
在\cref{ch:b1:poly}中，数字展开的类比是
按 $(X - a)$ 的幂展开。
\end{solution}
